\section{Transportes nonádicos y conservación de la memoria}
\label{sec:nonadic}

La reconstrucción de la sección~\ref{sec:memory} proporciona un enlace
efectivo entre la memoria del transporte y la geometría del registro
dodecafásico. El registro generado por APP--TRIT--TPK determina sus
lecturas de incidencia; las dos compresiones conservan, en componentes
separadas, la representación terminal y sus complementos; la reunión de
las memorias con la carga restituye las doce coordenadas. La dirección
geométrica y el proyector dimensional se obtienen después de esa
restitución. El desarrollo siguiente establece el balance operatorio
que sostiene esta cadena y su comportamiento bajo refinamiento,
realización espectral y cambio de métrica
\cite{HMT-VIII,hmtX}.

\subsection{Calendario de nueve fases y descomposición ortogonal}

Sea \(\mathcal H\) un espacio de Hilbert complejo y sea
\(C\in\mathcal B(\mathcal H)\). La suma ortogonal de nueve copias,
\(\mathcal H^{\oplus9}\), conserva por separado las nueve posiciones
del calendario. Definimos
\begin{equation}
 \begin{split}
 S_C(v_0,\ldots,v_8)&=(Cv_8,v_0,\ldots,v_7),\\
 J:\mathcal H\longrightarrow\mathcal H^{\oplus9},
 \qquad Jv&=\frac13(v,\ldots,v).
 \end{split}
 \label{nt:calendar}
\end{equation}
La normalización \(1/3=9^{-1/2}\) da \(J^*J=I\). El operador
\(P=JJ^*\) es la proyección ortogonal sobre la sección uniforme;
su acción sustituye las nueve componentes por su media. Cada
componente atraviesa una vez la unión que contiene \(C\) durante una
vuelta, de modo que
\begin{equation}
 S_C^9=I_9\otimes C.
 \label{nt:full-return}
\end{equation}
El retorno de la posición del calendario realiza así una acción de
\(C\) sobre la fibra de memoria.

\begin{definition}
La compresión uniforme del transporte y su complemento son
\begin{equation}
 T_C=J^*S_CJ=\frac{8I+C}{9},
 \qquad
 \eta_C=(I-P)S_CJ.
 \label{nt:compression}
\end{equation}
La primera aplicación tiene codominio \(\mathcal H\); la segunda
tiene codominio \(J(\mathcal H)^\perp\subset\mathcal H^{\oplus9}\).
\end{definition}

Para \(z=(C-I)v\), el complemento se escribe
\begin{equation}
 \eta_Cv=\frac1{27}(8z,-z,-z,-z,-z,-z,-z,-z,-z).
 \label{nt:complement}
\end{equation}
La suma de sus componentes es cero. Los coeficientes de la
descomposición proceden, por tanto, del promedio de ocho
continuaciones idénticas y una continuación transportada.

\begin{proposition}[Balance de la compresión y de la unión de memoria]
Para todo \(C\in\mathcal B(\mathcal H)\),
\begin{align}
 \eta_C^*\eta_C
   &=\frac8{81}(I-C)^*(I-C),\label{nt:eta-gram}\\
 I-T_C^*T_C
   &=\eta_C^*\eta_C+\frac19(I-C^*C).\label{nt:general-balance}
\end{align}
Si \(C^*C=I\), la aplicación
\(v\mapsto(T_Cv,\eta_Cv)\) es isométrica.
\end{proposition}

\begin{proof}
La expresión~\eqref{nt:complement} da
\[
 \|\eta_Cv\|^2
 =\frac{8^2+8}{27^2}\|(C-I)v\|^2
 =\frac8{81}\|(C-I)v\|^2.
\]
La polarización prueba~\eqref{nt:eta-gram}. Por otra parte,
\[
 T_C^*T_C
 =\frac{64I+8C+8C^*+C^*C}{81}.
\]
Sumando \eqref{nt:eta-gram} se obtiene
\[
 T_C^*T_C+\eta_C^*\eta_C=\frac{8I+C^*C}{9}.
\]
Restar este resultado de \(I\) da~\eqref{nt:general-balance}.
Cuando \(C^*C=I\), la suma de los dos productos de Gram es la
identidad, lo que conserva todos los productos interiores.
\end{proof}

El término \(\eta_C^*\eta_C\) mide la dispersión entre las nueve
componentes. El término \((I-C^*C)/9\) registra, por separado,
la variación de la norma en la unión de memoria. La realización
unitaria sitúa toda la variación visible en el primer término y
conserva la norma en la suma ortogonal completa.

\begin{proposition}[Conservación de un observable]
Sean \(C\) unitario y \(A=A^*\in\mathcal B(\mathcal H)\) tales que
\(C^*AC=A\). Entonces
\begin{equation}
 A=T_C^*AT_C+
   \eta_C^*(I_9\otimes A)\eta_C
  =T_C^*AT_C+\frac8{81}(I-C)^*A(I-C).
 \label{nt:observable-balance}
\end{equation}
Para \(A\succeq0\), ambos términos de la suma son positivos.
\end{proposition}

\begin{proof}
La fórmula del complemento, aplicada al producto sesquilineal
definido por \(A\), prueba la segunda igualdad. La expansión de
los dos términos da
\[
 \frac1{81}\bigl(
 64A+8AC+8C^*A+C^*AC
 +8A-8AC-8C^*A+8C^*AC\bigr)
 =\frac{8A+C^*AC}{9}=A.
\]
La positividad se conserva bajo una aplicación de la forma
\(B\mapsto L^*BL\).
\end{proof}

\subsection{Iteración de la compresión y conservación del sector fijo}

Fijamos en esta subsección un operador unitario \(C\) y escribimos
\(T=T_C\), \(\eta=\eta_C\). El transporte \(S_C^N\) mantiene todas
las fases durante \(N\) pasos. La potencia \(T^N\), en cambio,
aplica sucesivamente la compresión uniforme; en cada etapa
\(\eta T^jv\) conserva el complemento de la representación anterior.
La distinción queda expresada por los siguientes operadores.

\begin{proposition}[Registro isométrico a horizonte finito]
Para todo entero \(N\geq1\), la aplicación
\begin{equation}
 V_Nv=(T^Nv,\eta v,\eta Tv,\ldots,\eta T^{N-1}v)
 \label{nt:finite-record}
\end{equation}
de \(\mathcal H\) a
\(\mathcal H\oplus(\mathcal H^{\oplus9})^{\oplus N}\)
es isométrica. Si \(A\) satisface las hipótesis de
\eqref{nt:observable-balance}, entonces
\begin{equation}
 A=T^{*N}AT^N+
 \sum_{j=0}^{N-1}T^{*j}\eta^*(I_9\otimes A)\eta T^j.
 \label{nt:telescopy}
\end{equation}
\end{proposition}

\begin{proof}
Multiplicar~\eqref{nt:observable-balance} por \(T^{*j}\) y
\(T^j\) da
\[
 T^{*j}AT^j-T^{*(j+1)}AT^{j+1}
 =T^{*j}\eta^*(I_9\otimes A)\eta T^j.
\]
La suma telescópica prueba~\eqref{nt:telescopy}. Para \(A=I\),
el miembro derecho es el producto de Gram de \(V_N\).
El paso de \(N\) a \(N+1\) sustituye únicamente \(T^Nv\) por
\((T^{N+1}v,\eta T^Nv)\), conservando todas las componentes
de memoria previamente producidas.
\end{proof}

\begin{theorem}[Límite de la compresión y memoria ilimitada]
Sea \(P_{\mathrm{fix}}\) la proyección ortogonal sobre
\(\ker(I-C)\). Entonces
\begin{align}
 \operatorname*{s\!-\!lim}_{N\to\infty}T^N
   &=P_{\mathrm{fix}},\label{nt:strong-limit}\\
 I&=P_{\mathrm{fix}}+
 \sum_{j=0}^{\infty}T^{*j}\eta^*\eta T^j.\label{nt:infinite-record}
\end{align}
La serie converge fuertemente. Para todo observable acotado
autoadjunto que satisfaga \(C^*AC=A\), se conserva además
\begin{equation}
 A=P_{\mathrm{fix}}AP_{\mathrm{fix}}+
 \sum_{j=0}^{\infty}
 T^{*j}\eta^*(I_9\otimes A)\eta T^j.
 \label{nt:infinite-observable}
\end{equation}
\end{theorem}

\begin{proof}
El teorema espectral del operador unitario realiza \(C\) como
la función coordenada \(z\) sobre la circunferencia unidad.
La compresión corresponde a \(t(z)=(8+z)/9\), y
\begin{equation}
 1-|t(z)|^2=\frac8{81}|1-z|^2,
 \qquad |z|=1.
 \label{nt:circular-defect}
\end{equation}
Así, \(t(z)^N\) converge a cero en cada \(z\ne1\) y conserva
el valor uno en \(z=1\). La cota \(|t(z)^N|\leq1\) permite
aplicar convergencia dominada a cada medida espectral vectorial.
Se obtiene~\eqref{nt:strong-limit} y, por el mismo argumento,
\(T^{*N}\to P_{\mathrm{fix}}\) fuertemente. La acotación de
\(A\), junto con \(\|T^N\|\leq1\), da
\(T^{*N}AT^N\to P_{\mathrm{fix}}AP_{\mathrm{fix}}\).
El límite de~\eqref{nt:telescopy} prueba las dos identidades.
\end{proof}

El registro
\[
 V_\infty v=(P_{\mathrm{fix}}v,\eta v,\eta Tv,\eta T^2v,\ldots)
\]
es, por consiguiente, una isometría. El sector fijo conserva la
parte que persiste en la representación terminal; las demás
componentes permanecen en la memoria de los complementos.

Para el desplazamiento bilateral
\(C e_n=e_{n+1}\) en \(\ell^2(\mathbb Z)\), una sucesión fija
es constante. La condición de cuadrado-sumabilidad fuerza su
anulación, por lo que \(P_{\mathrm{fix}}=0\). Toda la norma
queda entonces representada en los complementos. El espectro de
este desplazamiento es la circunferencia completa, y el máximo
espectral de \(|t(z)|^N\) es uno; de ahí
\(\|T^N\|=1\) para cada \(N\). La convergencia fuerte conserva
este comportamiento no uniforme.

Para una permutación cíclica de \(d\) posiciones, el sector fijo
es la recta de los vectores uniformes. Si la permutación tiene
varias órbitas, cada órbita aporta su propia dirección uniforme.
Esta dimensión determina la parte terminal que debe conservarse.

\begin{proposition}[Compresión de un transporte de orden tres]
Sea \(C\) unitario y \(C^3=I\). Entonces
\begin{equation}
 P_{\mathrm{fix}}=\frac{I+C+C^2}{3},
 \qquad
 T_C^*T_C=P_{\mathrm{fix}}+
 \frac{19}{27}(I-P_{\mathrm{fix}}).
 \label{nt:order-three}
\end{equation}
\end{proposition}

\begin{proof}
El promedio de las tres potencias es el proyector ortogonal
sobre los vectores fijos. Como \(C^*=C^2\),
\[
 T_C^*T_C=\frac{65I+8(C+C^2)}{81}
 =\frac{57I+24P_{\mathrm{fix}}}{81},
\]
que coincide con~\eqref{nt:order-three}.
\end{proof}

En la realización de doce posiciones de
\(\mathbb R^{12}\), \((Sx)_i=x_{i+1}\) con índices módulo
doce, \(S^3\) tiene orden cuatro y \(S^4\) tiene orden tres.
Las cuatro órbitas de \(S^4\) dan
\(\operatorname{rank}P_{\mathrm{fix}}=4\).
La realización excepcional \(g_c\) del artículo~X satisface
\(g_c^3=I\) y \(\ker(I-g_c)=0\) en su propio portador;
en ella~\eqref{nt:order-three} se reduce a
\(T_{g_c}^*T_{g_c}=19I/27\)
\cite{hmtX}. El coeficiente y el sector fijo se transportan
conjuntamente con el dominio.

La contracción radial \(M=qC\), \(0<q<1\), posee otro registro
exacto:
\begin{equation}
 L_Nv=\bigl(\sqrt{1-q^2}\,v,
 \sqrt{1-q^2}\,qCv,\ldots,
 \sqrt{1-q^2}\,q^{N-1}C^{N-1}v,q^NC^Nv\bigr).
 \label{nt:radial-record}
\end{equation}
Su producto de Gram es la identidad porque
\((1-q^2)\sum_{j=0}^{N-1}q^{2j}+q^{2N}=1\).
El terminal tiene norma \(q^N\|v\|\), y el registro infinito
conserva el producto interior. El factor \(q=5/9\) del modo
diferencial central del artículo~VIII se aplica a una vuelta
completa; su reparto uniforme entre nueve pasos utiliza
\(q^{1/9}S_C\), cuya novena potencia es
\(q(I_9\otimes C)\) \cite{HMT-VIII}.

\subsection{Reducción del registro y conservación de la traza}

La dirección fija del complemento en las nueve componentes
permite una realización mínima en dos sumandos. Definamos
\[
 D_C=\frac{\sqrt8}{9}(I-C).
\]
Las identidades anteriores dan
\(T_C^*T_C+D_C^*D_C=I\) para \(C\) unitario.
El cambio de signo respecto de la expresión de
\(\eta_C\) corresponde a una orientación de la componente
complementaria y conserva su producto de Gram.

\begin{proposition}[Canal positivo de la compresión con memoria]
Para \(C\) unitario, la aplicación
\begin{equation}
 \mathcal Q_C(X)=T_CXT_C^*+D_CXD_C^*
              =\frac{8X+CXC^*}{9}
 \label{nt:channel}
\end{equation}
es completamente positiva en \(\mathcal B(\mathcal H)\) y
conserva la identidad. Su restricción a los operadores de
clase traza conserva la traza.
\end{proposition}

\begin{proof}
Al expandir los dos términos de la primera expresión,
los productos \(CX\) y \(XC^*\) se cancelan. Los coeficientes
restantes son \(72/81=8/9\) para \(X\) y \(9/81=1/9\)
para \(CXC^*\), lo que prueba~\eqref{nt:channel}.
Para cualquier entero \(m\geq1\), la ampliación a matrices
de orden \(m\) tiene la forma
\[
 X\longmapsto
 (I_m\otimes T_C)X(I_m\otimes T_C)^*
 +(I_m\otimes D_C)X(I_m\otimes D_C)^*.
\]
Cada sumando conserva positividad; por ello la aplicación es
completamente positiva. La segunda expresión de
\eqref{nt:channel} da \(\mathcal Q_C(I)=I\).
Si \(X\) es de clase traza, la invariancia de la traza bajo
conjugación unitaria implica
\(\operatorname{Tr}\mathcal Q_C(X)=\operatorname{Tr}X\).
\end{proof}

La aplicación \eqref{nt:channel} se obtiene igualmente
preparando la isometría
\(v\mapsto T_Cv\otimes e_0+D_Cv\otimes e_1\) y efectuando
la traza parcial sobre \(\mathbb C^2\).
La conservación del vector completo pertenece a la isometría;
la conservación de la traza del estado reducido pertenece a
\(\mathcal Q_C\). Ambas afirmaciones proceden del mismo
balance ortogonal.

En la prolongación nonádica de álgebras finitas aparece una
conservación complementaria. Para
\(\mathcal A_d=M_{9^d}(\mathbb C)\), con
\(\tau_d(a)=9^{-d}\operatorname{Tr}a\), las aplicaciones
\[
 \iota_d(a)=a\otimes I_9,\qquad
 \jmath_{d,j}(a)=a\otimes p_j,
 \quad p_j=|e_j\rangle\langle e_j|
\]
satisfacen
\begin{align}
 \tau_{d+1}(\iota_d(a))&=\tau_d(a),
 &\iota_d(I)&=I,\label{nt:unital-trace}\\
 \tau_{d+1}(\jmath_{d,j}(a))&=\frac{\tau_d(a)}9,
 &\jmath_{d,j}(I)&=I\otimes p_j.\label{nt:corner-trace}
\end{align}
La demostración usa
\(\operatorname{Tr}(a\otimes b)=
\operatorname{Tr}(a)\operatorname{Tr}(b)\),
\(\operatorname{Tr}I_9=9\) y \(\operatorname{Tr}p_j=1\).
La primera inclusión conserva toda la fibra y la unidad;
la segunda conserva una continuación marcada y su soporte.

Esta distinción también se manifiesta en el árbol de
historias enriquecidas. Sean \(\mathcal H_n\) los prefijos
supervivientes de longitud \(n\), y sea \(d(h)\geq1\) el número de
emisiones salientes conservadas de \(h\). Cada emisión con multiplicidad
se representa por sucesores marcados distintos; la suma siguiente recorre
ese conjunto de \(d(h)\) sucesores.
La ley de pesos
\[
 \mu_{n+1}(h\epsilon)=\mu_n(h)\,p_h(\epsilon),
 \qquad p_h(\epsilon)=d(h)^{-1}
\]
conserva
\(\sum_{\epsilon}\mu_{n+1}(h\epsilon)=\mu_n(h)\).
Por tanto, en
\(\mathscr L_n=L^2(\mathcal H_n,\mu_n)\), el pullback
\[
 R_nf(h\epsilon)=f(h)
\]
es isométrico y conserva la función unidad. Su adjunto es
\[
 E_ng(h)=\sum_{\epsilon}p_h(\epsilon)g(h\epsilon),
 \qquad E_nR_n=I,\quad E_n1=1.
\]
En efecto, agrupar la suma de \(\|R_nf\|^2\) por predecesor
produce \(\sum_h\mu_n(h)|f(h)|^2\); el mismo agrupamiento
en el producto interior determina \(E_n\).
La memoria y las multiplicidades quedan presentes en el
dominio de las historias. Estos mapas constituyen la
realización funcional del refinamiento descrito en el
artículo~X \cite{hmtX}.

\subsection{Coordenatización espectral y reciprocidad del radio}

El transporte anterior admite realizaciones espectrales
con dominios propios. La coordenatización de Möbius de VIII actúa,
en coordenadas finitas, por
\begin{equation}
 \tau(z)=\frac{z-1}{z},
 \qquad z\in\mathbb C\setminus\{0\}.
 \label{nt:mobius}
\end{equation}
La inversa es \(z=(1-\tau)^{-1}\), para \(\tau\ne1\).
La extensión a la esfera de Riemann envía \(0\) a
\(\infty\) e \(\infty\) a \(1\).

\begin{proposition}[Espejo espectral y radio recíproco]
Sea \(z^\sharp=1-\overline z\). Para
\(z\notin\{0,1\}\),
\begin{equation}
 \tau(z^\sharp)=\frac1{\overline{\tau(z)}}.
 \label{nt:mirror}
\end{equation}
Si \(z=\sigma+it\), entonces
\begin{equation}
 |\tau(z)|^2
 =\frac{(\sigma-1)^2+t^2}{\sigma^2+t^2};
 \qquad
 |\tau(z)|=1\ \Longleftrightarrow\ \sigma=\frac12.
 \label{nt:critical-circle}
\end{equation}
\end{proposition}

\begin{proof}
La sustitución directa da
\[
 \tau(1-\overline z)
 =\frac{-\overline z}{1-\overline z}
 =\frac{\overline z}{\overline z-1}
 =\frac1{\overline{(z-1)/z}}.
\]
El cociente de módulos cuadrados produce
\eqref{nt:critical-circle}; la igualdad entre numerador y
denominador equivale a \(-2\sigma+1=0\).
\end{proof}

Escribiendo \(\tau(z)=r e^{i\theta}\), con \(r>0\), el
espejo preserva \(\theta\) y transforma \(r\) en \(r^{-1}\).
Su conjunto fijo en esta coordenatización es la circunferencia unidad.
La conjugación \(z\mapsto\overline z\), por su parte,
produce \(\tau\mapsto\overline\tau\) y cambia la orientación
angular. Las dos involuciones conservan, respectivamente,
la coordenada angular y la coordenada radial.

Sobre \(z=1/2+i\gamma\), con \(\gamma\in\mathbb R\),
\begin{equation}
 \tau_\gamma=\frac{\gamma+i/2}{\gamma-i/2},
 \qquad
 \gamma=\frac{i}{2}\frac{\tau_\gamma+1}{\tau_\gamma-1}
       =\frac12\cot\frac{\theta}{2},
 \quad \tau_\gamma=e^{i\theta}.
 \label{nt:spectral-phase}
\end{equation}
La compactificación conserva la altura mediante la inversa;
el punto \(\tau=1\) corresponde al infinito.

Sea \(\mathcal D=C_c^\infty(\mathbb R;\mathbb C)\), con su topología
usual de funciones de prueba. Para la realización positiva de la forma
de Weil se utilizan expresamente su positividad semidefinida, su
invariancia por traslación y su continuidad en esa topología, en el
dominio demostrado por el lector espectral correspondiente. Si
\(\mathscr W(f,g)=\langle Ef,Eg\rangle\), su núcleo es
\(\mathcal N=\{f\in\mathcal D:\mathscr W(f,f)=0\}\), y se define
\[
 \mathcal H_{\mathscr W}
 =\overline{\mathcal D/\mathcal N}.
\]
Las traslaciones inducen un grupo unitario fuertemente
continuo \(U_t[f]=[f(\,\cdot-t)]\): la invariancia preserva la norma,
y la continuidad de traslaciones en \(\mathcal D\) y de la forma
prueba la continuidad fuerte sobre un subespacio denso y después en
toda la completación. Con la convención
\(U_t=e^{itH}\), su generador autoadjunto actúa sobre
las clases de pruebas por \(H[f]=[if']\). El Cayley de
esta realización es el operador acotado
\begin{equation}
 C_{\mathscr W}
 =(H+iI/2)(H-iI/2)^{-1}
 =I+i(H-iI/2)^{-1}.
 \label{nt:cayley-hilbert}
\end{equation}
El cálculo funcional lo hace unitario, pues el cociente
en~\eqref{nt:spectral-phase} tiene módulo uno para
\(\gamma\) real. La identidad de compresión se aplica
a \(C_{\mathscr W}\) en \(\mathcal H_{\mathscr W}\).
El espacio, el dominio del generador y la positividad
que permite su construcción son las premisas de esta
realización espectral \cite{HMT-VIII}.

La memoria bilateral actúa, en cambio, en
\(\ell^2(\mathbb Z)\), y las permutaciones \(S^3,S^4\)
actúan en \(\mathbb R^{12}\). Cada realización determina
su espectro, sus sectores fijos y la acción de una vuelta.
Los respectivos mapas de transporte especifican las
composiciones entre estos espacios.

\subsection{Cayley sobre pruebas y conservación de la forma aritmética}

La misma representación aritmética admite un lector de Cayley
definido directamente sobre funciones de prueba. Fijemos
\(a>b>1/2\) y
\[
 \mathscr E_b=
 \left\{f\in C^\infty(\mathbb R):
   \sup_x e^{b|x|}|f^{(j)}(x)|<\infty
   \text{ para todo }j\geq0\right\}.
\]
Las funciones suaves de soporte compacto pertenecen a
\(\mathscr E_b\). Definimos
\begin{align}
 (R_a^+f)(x)&=\int_0^\infty e^{-at}f(x+t)\,dt,
 &(R_a^-f)(x)&=\int_0^\infty e^{-at}f(x-t)\,dt,\label{nt:resolvents}\\
 C_af&=f-2aR_a^+f,
 &C_a^{-1}f&=f-2aR_a^-f.\label{nt:test-cayley}
\end{align}
Cada derivada de \(R_a^\pm f\) satisface la cota del
correspondiente seminorma de \(f\), multiplicada por
\((a-b)^{-1}\). Esto sigue de
\(e^{b|x|}|f(x\pm t)|\leq M_0e^{bt}\).
Ambas aplicaciones preservan \(\mathscr E_b\).

Con la transformada
\(\widehat f(\xi)=\int_{\mathbb R}f(x)e^{-i\xi x}\,dx\),
la aplicación de Fubini da
\begin{equation}
 \widehat{C_af}(\xi)
 =c_a(\xi)\widehat f(\xi),\qquad
 c_a(\xi)=\frac{\xi-ia}{\xi+ia}.
 \label{nt:cayley-multiplier}
\end{equation}
El multiplicador de la segunda fórmula
de~\eqref{nt:test-cayley} es \(c_a^{-1}\), lo que prueba
la inversa. Para \(\xi\in\mathbb R\),
\(|c_a(\xi)|=1\); Plancherel proporciona su extensión
unitaria a \(L^2(\mathbb R)\).

Para precisar la conservación aritmética, sean
\[
 \mathcal C_{f,g}(u)
 =\int_{\mathbb R}f(x)\overline{g(x-u)}\,dx,
 \qquad
 P_f^\pm=\int_{\mathbb R}f(x)e^{\pm x/2}\,dx.
\]
Las longitudes \(\ell_{p,k}=k\log p\) y los pesos
\(w_{p,k}=(\log p)p^{-k/2}\) son las representaciones
primo--potencia del lector aritmético. La forma completa
del artículo~VIII se escribe
\begin{equation}
 \begin{split}
 \mathscr W(f,g)={}&c_\Gamma\mathcal C_{f,g}(0)\\
 &+\int_0^\infty\frac{e^{-u/2}}{1-e^{-2u}}
 \bigl[2\mathcal C_{f,g}(0)
       -\mathcal C_{f,g}(u)-\mathcal C_{f,g}(-u)\bigr]\,du\\
 &-\sum_{p,k}w_{p,k}
   \bigl[\mathcal C_{f,g}(\ell_{p,k})
        +\mathcal C_{f,g}(-\ell_{p,k})\bigr]
 +P_f^+\overline{P_g^-}+P_f^-\overline{P_g^+},
 \end{split}
 \label{nt:weil-complete}
\end{equation}
donde \(c_\Gamma=\psi(1/4)-\log\pi\), y
\(\psi=\Gamma'/\Gamma\) es la derivada logarítmica de la
función gamma. Estas funciones y constantes intervienen
en la realización analítica posterior del lector
aritmético del corpus.

La cota
\[
 |\mathcal C_{f,g}(u)|
 \leq M_fM_g\bigl(|u|+b^{-1}\bigr)e^{-b|u|}
\]
controla la suma prima por una serie convergente del tipo
\(\sum_{n\geq2}(\log n)(1+\log n)n^{-b-1/2}\).
La diferencia simétrica de correlaciones es \(O(u^2)\)
en el origen, mientras que el peso gamma es \(O(u^{-1})\).
En infinito ese peso decrece exponencialmente.
El margen \(b>1/2\) asegura asimismo la convergencia de
los dos lectores polares. Así queda fijado el dominio
de todos los términos de~\eqref{nt:weil-complete}.

\begin{proposition}[Invariancia aritmética y balance nonádico]
Para \(f,g\in\mathscr E_b\),
\begin{equation}
 \mathscr W(C_af,C_ag)=\mathscr W(f,g).
 \label{nt:weil-invariance}
\end{equation}
Con \(\mathcal T_a=(8I+C_a)/9\), se cumple
\begin{equation}
 \mathscr W(f,g)
 =\mathscr W(\mathcal T_af,\mathcal T_ag)
 +\frac8{81}\mathscr W((I-C_a)f,(I-C_a)g).
 \label{nt:weil-balance}
\end{equation}
\end{proposition}

\begin{proof}
El multiplicador \(c_a\) es unitario en la recta real y
conmuta con las traslaciones. En consecuencia,
\(\mathcal C_{C_af,C_ag}(u)=\mathcal C_{f,g}(u)\)
para todo \(u\). Los lectores polares se transforman
por Fubini según
\[
 P_{C_af}^+
 =-\frac{a-1/2}{a+1/2}P_f^+,\qquad
 P_{C_af}^-
 =-\frac{a+1/2}{a-1/2}P_f^-.
\]
Sus factores son reales y recíprocos, de modo que
cada producto polar cruzado permanece inalterado.
La reunión de los términos de
\eqref{nt:weil-complete} prueba~\eqref{nt:weil-invariance}.
La expansión sesquilineal del segundo miembro de
\eqref{nt:weil-balance} cancela los términos mixtos y
produce
\(\bigl(8\mathscr W(f,g)+
\mathscr W(C_af,C_ag)\bigr)/9\), igual a
\(\mathscr W(f,g)\).
\end{proof}

La iteración finita conserva la forma completa:
\[
 \mathscr W(f,g)
 =\mathscr W(\mathcal T_a^Nf,\mathcal T_a^Ng)
 +\frac8{81}\sum_{j=0}^{N-1}
 \mathscr W((I-C_a)\mathcal T_a^jf,
            (I-C_a)\mathcal T_a^jg).
\]
Esta identidad sesquilineal conserva simultáneamente
gamma, primos y términos polares. Su evaluación en una
realización positiva permite interpretarla como reparto
de energía; su validez algebraica precede a esa
interpretación. El margen \(a>b>1/2\) pertenece al
dominio exponencial de pruebas de esta construcción.
El Cayley~\eqref{nt:cayley-hilbert} con parámetro \(1/2\)
pertenece al dominio hilbertiano de su generador.

\subsection{Covariancia métrica del flujo dodecafásico}

La reconstrucción de la sección~\ref{sec:memory} determina
\(u_K\in\mathbb R^{12}\), con
\(\mathbf1^{\mathsf T}u_K=0\) y \(u_K\ne0\).
Escribimos \(A_K=\operatorname{diag}(u_K)\), con la dirección y la
normalización declaradas en \eqref{eq:u-main}, y
\[
 g_s=e^{sA_K}\in\operatorname{GL}(12,\mathbb R).
\]
La traza nula de \(A_K\) implica \(\det g_s=1\).
Cada \(g_s\) transporta los pesos y las coordenadas
geométricas según la realización correspondiente.
Para conservar también el balance de memoria entre
coordenatizaciones hay que transportar el producto interior.

\begin{proposition}[Transporte métrico de compresión y complemento]
Sean \(V\) un espacio hermítico finito, \(C_0\) unitario
y \(g:V\to V\) invertible. Escribimos
\(g^{-*}=(g^{-1})^*\) y definimos
\begin{equation}
 C_g=gC_0g^{-1},\qquad
 G_g=g^{-*}g^{-1},\qquad
 \langle v,w\rangle_g=\langle G_gv,w\rangle.
 \label{nt:metric-chart}
\end{equation}
Entonces \(C_g^*G_gC_g=G_g\), y
\begin{align}
 T_{C_g}g&=gT_{C_0},\label{nt:metric-naturality}\\
 \eta_{C_g}g&=g^{\oplus9}\eta_{C_0},\label{nt:eta-naturality}\\
 G_g&=T_{C_g}^*G_gT_{C_g}
 +\frac8{81}(I-C_g)^*G_g(I-C_g).\label{nt:metric-balance}
\end{align}
\end{proposition}

\begin{proof}
La positividad de \(G_g\) se obtiene de
\(\langle G_gv,v\rangle=\|g^{-1}v\|^2\).
Sustituyendo~\eqref{nt:metric-chart},
\[
 C_g^*G_gC_g
 =g^{-*}C_0^*g^*g^{-*}g^{-1}gC_0g^{-1}
 =g^{-*}C_0^*C_0g^{-1}=G_g.
\]
La identidad \(C_gg=gC_0\) da
\eqref{nt:metric-naturality} al sumar \(8g\) y dividir
por nueve. También da
\((C_g-I)g=g(C_0-I)\); su sustitución en
\eqref{nt:complement} prueba~\eqref{nt:eta-naturality}.
Finalmente, la expansión de
\eqref{nt:observable-balance}, con \(A=G_g\) y
\(C_g^*G_gC_g=G_g\), prueba
\eqref{nt:metric-balance}.
\end{proof}

Aplicando la proposición a \(g=g_s\) y a los operadores
\(C_0=S^3,S^4\) de la realización dodecafásica, las
compresiones, las memorias y sus productos interiores
se transportan por cuadrados conmutativos explícitos.
La identificación \(g_s:(V,\langle\cdot,\cdot\rangle)
\to(V,\langle\cdot,\cdot\rangle_{g_s})\) es isométrica.
Por tanto, la fórmula transporta también los sectores
fijos y los registros a toda profundidad:
\[
 P_{\mathrm{fix},s}
 =g_sP_{\mathrm{fix},0}g_s^{-1}.
\]
El proyector del miembro derecho es ortogonal respecto
de \(G_{g_s}\).

La condición de volumen y la condición métrica tienen
efectos diferenciados. Para la matriz racional
\[
 g=\begin{pmatrix}2&0\\0&1/2\end{pmatrix},
 \qquad \det g=1,
\]
la sustitución directa en la compresión da
\[
 \frac{8I+g}{9}
 =\begin{pmatrix}10/9&0\\0&17/18\end{pmatrix},
 \qquad
 T_g^*T_g+\frac8{81}(I-g)^*(I-g)
 =\frac{8I+g^*g}{9}.
\]
La primera coordenada se expande en la métrica
euclídea. El balance~\eqref{nt:metric-balance}
conserva en cambio la unidad métrica cuando se
transporta un operador unitario \(C_0\) junto con
su producto interior. El determinante controla
el volumen; \(G_g\) controla las normas y los
productos cruzados.

El flujo satisface \(g_{-s}=g_s^{-1}\).
La reciprocidad del radio de
\eqref{nt:mirror} pertenece a la coordenatización espectral,
y la inversión de \(g_s\) pertenece a la
realización dimensional del registro. El enlace
constructivo utilizado aquí conserva las
aplicaciones explícitas de memoria a \(K\) y de
\(K\) a su dirección geométrica, seguido por
el transporte métrico de cada realización.
Esta composición mantiene, en un mismo desarrollo,
la reconstrucción exacta, el refinamiento positivo,
el registro de fronteras y la conservación
de los productos interiores.
